% Part: model-theory
% Chapter: lindstrom
% Section: ls-property

\documentclass[../../../include/open-logic-section]{subfiles}

\begin{document}

\olfileid{mod}{lin}{lsp}

\olsection{ویژگی‌های فشردگی و لوونهایم--اسکولم}

اکنون تعمیم‌های طبیعی فشردگی و لوونهایم--اسکولم را به منطق‌های
انتزاعی ارائه می‌کنیم.

\begin{defn}
منطق انتزاعی~$\tuple{L, \models_L}$
\emph{ویژگی فشردگی} دارد اگر هر مجموعهٔ~$\Gamma$ از !!{sentence}s
$L(\Lang{L})$، هرگاه هر $\Gamma_0 \subseteq \Gamma$ متناهی ارضاپذیر
باشد، ارضاپذیر باشد.
\end{defn}

\begin{defn}
$\tuple{L, \models_L}$
\emph{ویژگی لوونهایم--اسکولم نزولی} دارد اگر هر~$\Gamma$ ارضاپذیر
!!a{enumerable} مدل داشته باشد.
\end{defn}


مفهوم همریختی جزئی از
\olref[bas][pis]{defn:partialisom} کاملاً «جبری» است؛ یعنی بدون ارجاع
به !!{sentence}s زبان و تنها برحسب محمول‌ها و ثابت‌هایی تعریف شده است
که !!{language}~$\Lang{L}$ در اختیار !!{structure}s می‌گذارد. پس این
مفهوم برای منطق‌های انتزاعی نیز کاربرد دارد. دربارهٔ منطق مرتبهٔ اول،
از
\olref[bas][pis]{thm:p-isom2} می‌دانیم اگر دو !!{structure}s تا حدی
همریخت باشند، هم‌ارز مقدماتی‌اند. آن اثبات به منطق‌های انتزاعی منتقل
نمی‌شود، زیرا ممکن است استقرا روی !!{formula}s برای
$!E \in L(\Lang{L})$ دلخواه در دسترس نباشد؛ بااین‌حال، اگر ویژگی
لوونهایم--اسکولم برقرار باشد، خود قضیه همچنان درست است.

\begin{thm}
\ollabel{thm:abstract-p-isom}
فرض کنید $\tuple{L, \models_L}$ منطقی نرمال با ویژگی
لوونهایم--اسکولم باشد. آنگاه هر دو !!{structure}s که تا حدی همریخت
باشند، در~$\tuple{L,\models_L}$ هم‌ارز مقدماتی‌اند.
\end{thm}

\begin{proof}
فرض کنید $\Struct{M} \simeq_p \Struct{N}$، اما برای یک~$!E$،
$\Struct{M} \models_L !E$ درحالی‌که
$\Struct{N} \not\models_L !E$. بنا بر ویژگی همریختی می‌توانیم فرض کنیم
$\Domain{M}$ و~$\Domain{N}$ جدا از هم‌اند، و بنا بر ویژگی بسط
می‌توانیم فرض کنیم $!E \in L(\Lang{L})$ برای !!{language} متناهی
$\Lang{L}$. بگذارید $\mathcal{I}$ مجموعه‌ای از همریختی‌های جزئی میان
$\Struct{M}$ و~$\Struct{N}$ باشد؛ بی‌کاستن از کلیت، همچنین فرض کنید
اگر $p \in \mathcal{I}$ و $q \subseteq p$، آنگاه
$q \in \mathcal{I}$.

$\Domain{M}^{<\omega}$ مجموعهٔ دنباله‌های متناهی از !!{element}s
$\Domain{M}$ است. بگذارید $S$ رابطه‌ای سه‌موضعی روی
$\Domain{M}^{<\omega}$ باشد که الحاق را بازنمایی می‌کند؛ یعنی اگر
$\mathbf{a}, \mathbf{b}, \mathbf{c} \in \Domain{M}^{<\omega}$،
آنگاه $S(\mathbf{a}, \mathbf{b}, \mathbf{c})$ هرگاه و تنها هرگاه
برقرار است که $\mathbf{c}$ حاصل الحاق~$\mathbf{a}$ و~$\mathbf{b}$
باشد. همچنین بگذارید $T$ رابطه‌ای سه‌موضعی باشد چنان‌که برای
$b \in \Domain{M}$ و
$\mathbf{a}, \mathbf{c} \in \Domain{M}^{<\omega}$،
$T(\mathbf{a}, b, \mathbf{c})$ هرگاه و تنها هرگاه برقرار باشد که
$\mathbf{a}=a_1,\dots,a_n$ و
$\mathbf{c}=a_1,\dots,a_n,b$. !!{predicate}s سه‌موضعی تازهٔ~$P$ و
$Q$ را برگزینید و !!{structure}~$\Struct{M}^*$ را با جهان
$\Domain{M} \cup \Domain{M}^{<\omega}$ بسازید، به‌گونه‌ای که
$\Struct{M}$ زیرساختار آن باشد و $P$ و~$Q$ به‌ترتیب با روابط
الحاق~$S$ و~$T$ تعبیر شوند (پس $\Struct{M}^*$ در !!{language}
$\Lang{L} \cup \{P,Q\}$ است).

$\Domain{N}^{<\omega}$، $S'$، $T'$، $P'$، $Q'$ و
$\Struct{N}^*$ را به‌طور مشابه تعریف کنید. چون بنا بر فرض
$\Struct{M} \simeq_p \Struct{N}$، رابطه‌ای چون~$I$ میان
$\Domain{M}^{<\omega}$ و $\Domain{N}^{<\omega}$ وجود دارد چنان‌که
$I(\mathbf{a},\mathbf{b})$ هرگاه و تنها هرگاه برقرار است که نگاشت
مؤلفه‌های~$\mathbf{a}$ به مؤلفه‌های متناظر~$\mathbf{b}$ همریختی جزئی
باشد و شرط رفت‌وبرگشتِ
\olref[bas][pis]{defn:partialisom} را برآورده کند. اکنون بگذارید
$\Struct{A}$ !!{structure}ای باشد که !!{domain} آن اجتماع
!!{domain}s دو ساختار~$\Struct{M}^*$ و~$\Struct{N}^*$ است، این دو را
به‌عنوان زیر!!{structure}s دربر می‌گیرد، و !!{language} آن یک !!{predicate}
دوموضعی افزودهٔ~$R$ دارد که با رابطهٔ~$I$ تعبیر می‌شود؛ زبان همچنین
!!{predicate}sی دارد که !!{domain}s
$\Domain{M}^*$ و~$\Domain{N}^*$ را مشخص می‌کنند.

\begin{figure}[h]
  \centering
  \begin{tikzpicture}[node distance=2cm, auto, thick, >=stealth']
    \draw [rounded corners] (0,0) -- (8,0) -- (8,4) -- (0,4) --  cycle;
    \draw (2,2) circle (0.5cm);
    \draw (2,2) circle (1.25cm);
    \draw (6,2) circle (0.5cm);
    \draw (6,2) circle (1.25cm);
    \path node at (0.75,3.5) {\large $\Struct{A}$};
    \path node at (2,2) {\large $\Struct{M}$};
    \path node at (6,2) {\large $\Struct{N}$};
    \path node at (3.5,1) {\large $\Struct{M}^*$};
    \path node at (7.5,1) {\large $\Struct{N}^*$};
    \node (Idom) at (2.8,2) {};
    \node (Irng) at (5.2,2) {};
    \draw[<->, bend left] (Idom) to node {\large $I$} (Irng) ;
  \end{tikzpicture}
  \caption{!!{structure}~$\Struct{A}$ با همریختی جزئی درونی.}
\end{figure}

مشاهدهٔ تعیین‌کننده این است که در !!{language} این !!{structure}،
یعنی~$\Struct{A}$، به کمک ویژگی نسبی‌سازی، !!{sentence}~$!D_1$ از
منطق انتزاعی~$L$ وجود دارد
که در~$\Struct{A}$ صادق است و می‌گوید
$\Struct{M} \models_L !E$ و $\Struct{N} \not\models_L !E$. همچنین
!!{sentence}
\emph{مرتبهٔ اول}~$!D_2$ در~$\Struct{A}$ صادق است که می‌گوید
$\Struct{M} \simeq_p \Struct{N}$ از راه همریختی جزئی~$I$. چون منطق
نرمال همهٔ جمله‌های مرتبهٔ اول را دربر می‌گیرد و نسبت به عطف بسته است،
$!D_1$ و~$!D_2$ را می‌توان با هم در~$L$ بیان کرد. بنا بر ویژگی
لوونهایم--اسکولم، این دو جمله مشترکاً در !!a{enumerable} مدل
$\Struct{A}_0$ صادق‌اند؛ مدلی که زیرساختارهای تا حدی همریخت
$\Struct{M}_0$ و~$\Struct{N}_0$ را دربر دارد، چنان‌که
$\Struct{M}_0 \models_L !E$ و
$\Struct{N}_0 \not\models_L !E$. اما !!{enumerable} !!{structure}s تا حدی
همریخت، بنا بر
\olref[bas][pis]{thm:p-isom1}، در واقع همریخت‌اند؛ و این با ویژگی
همریختی منطق‌های انتزاعی نرمال در تناقض است.
\end{proof}

\end{document}
